\section{Conclusion and Future Work}

We presented \textbf{ZhuLong-SE (Self-Evolution)}, which augments ZhuLong~\cite{liu2026zhulong} with task-level harness enhancements and governed cross-task evolution while keeping the backbone frozen. On the 158-task EDA-Eval-PyAether benchmark, the full task-level harness achieves 92.4\% pass@1 in a single full run and outperforms ZhuLong on all five backbone comparisons. On the 40-task instance-disjoint challenge test, self-evolution from the Bare configuration achieves 80.0\% mean pass@1, compared with 77.5\% for the human-crafted-experience baseline. The frozen-trace ablation identifies the effects of ground-truth anchoring, contradiction handling, and routing on reflection output, while leaving scope-boundary reflection and multi-batch disconfirmation unverified under its entry-count, single-round design. Future work will extend the evaluation to additional EDA tool flows and languages, including SKILL-based automation, and to larger task sets, while investigating adaptive gating thresholds that adjust to task difficulty.
% Its static checker targets deterministic API misuse and raises no flags on the 158 ground-truth solutions, a sanity check limited to this reference-code set. The self-evolution layer distills traces into a gated three-tier memory store, rejects degrading rounds through held-out validation, and removes full call sequences and task-specific details during retrieval. In our held-out EDA evaluation, this design achieves a modest mean improvement over a human-crafted-experience baseline.
% 我们提出了\textbf{ZhuLong-SE（Self-Evolution）}，在ZhuLong~\cite{liu2026zhulong}之上加入任务级框架增强和受治理的跨任务自进化，同时保持基座模型冻结。其静态检查器针对确定性API误用，在158个ground-truth解上未标记任何问题，此健全性检查仅限于该参考代码集。自进化层将轨迹蒸馏到受门控的三层记忆存储中，通过留出验证拒绝退化轮次，并在检索时移除完整调用序列和任务特定细节。在本文的留出EDA评估中，该设计相比人工经验基线取得了适度的平均提升。
% 在158任务基准上，完整任务级框架单次全量运行达到92.4\%，并在当前所有的五组基座比较中均优于ZhuLong。在40任务实例不相交挑战测试上，从Bare配置开始的自进化达到80.0\%平均pass@1，而人工经验基线为77.5\%。冻结轨迹消融识别了真值锚定、矛盾处理和路由对反思输出的影响，而范围边界反思和多批次证伪在当前条目计数、单轮设计下尚未验证。

% Limitations include the modest challenge test set (40 tasks) and evaluation on a single EDA domain; broader generalization remains to be explored.
% 局限性包括挑战测试集规模有限（40个任务）以及在单一EDA领域的评估；更广泛的泛化性仍有待探索。

% Beyond scaling, we plan to push along the RSI trajectory toward making the harness code itself an optimization target, following the Meta-Harness paradigm~\cite{lee2026metaharness}—our governance gates already provide the validation infrastructure required for safe harness self-modification. A key challenge is the evaluation bottleneck identified in the RSI literature: harness-level evolution requires more nuanced quality metrics beyond pass/fail to avoid reward hacking. We also plan to explore whether the governed-memory architecture can serve as a general-purpose self-evolution layer for other API-driven automation domains, such as simulation and verification scripting.
% 未来工作将把评估扩展到更多EDA工具流程和更大的任务集，并研究适应任务难度的自适应门控阈值。在规模扩展之外，我们计划沿着RSI轨迹推进，使Harness代码本身成为优化目标，遵循Meta-Harness范式~\cite{lee2026metaharness}——我们的治理门控已为安全的Harness自我修改提供了验证基础设施。一个关键挑战是RSI文献中指出的评估瓶颈：Harness级进化需要超越通过/失败的更精细质量度量，以避免奖励黑客。我们还计划探索受治理记忆架构是否可作为其他API驱动的自动化领域（如仿真和验证脚本编写）的通用自进化层。
